problem:  
Let \((M,g(t))_{t\ge0}\) be a smooth family of closed \(n\)-dimensional Riemannian manifolds evolving under the **normalized Ricci flow**
\[
\partial_t g = -2(\mathrm{Ric}(g) - \tfrac{R_\mathrm{avg}}{n} g), 
\]
and let \(V:\mathbb{R}\to\mathbb{R}\) be a smooth double‑well potential with minima at \(u_\pm\).  
For \(\varepsilon>0\), consider the **curvature‑coupled Allen–Cahn system**
\[
\partial_t u_\varepsilon = \Delta_{g(t)}u_\varepsilon - \frac{1}{\varepsilon^2}V'(u_\varepsilon) + \alpha R_{g(t)}u_\varepsilon.
\]
Define the **energy measure**
\[
\mu^t_\varepsilon := \left(\frac{\varepsilon}{2}|\nabla u_\varepsilon|_{g(t)}^2 + \frac{1}{\varepsilon}V(u_\varepsilon)\right)\!d\mu_{g(t)}.
\]
Assume a uniform energy bound \(\sup_{t\le T}\mu^t_\varepsilon(M)<C_T\), and that the metrics \(g(t)\) converge smoothly (up to diffeomorphism) to an Einstein metric \(g_\infty\).

**Open problem (Allen–Cahn–Ricci flow interface conjecture):**  
As \(\varepsilon\to0\), for each fixed \(T>0\), does the family \(\{\mu^t_\varepsilon\}_{t\in[0,T]}\) subconverge (in the sense of integral varifolds depending on \(t\)) to a limit family \(V_t\) representing a **generalized mean‑curvature flow** in the evolving metric \(g(t)\)?  
More precisely:

1. Can one prove that any such limit \(V_t\) satisfies a Brakke‑type inequality in \((M,g(t))\),
   \[
   \frac{d}{dt}\!\int\!\phi\, d\|V_t\| 
   \le -\!\int (\phi |H_t|^2 - \langle\nabla_{g(t)}\phi, H_t\rangle)\,d\|V_t\|
   +\mathcal{K}[\phi,g(t)], 
   \]
   where \(\mathcal{K}[\phi,g(t)]\) is an explicit curvature‑dependent correction arising from \(\partial_t g\) and \(\alpha R_{g(t)}\)?  

2. If this evolving‑metric Brakke inequality holds, does \(V_t\) converge as \(t\to\infty\) to a stationary varifold \(V_\infty\)?  Is \(V_\infty\) necessarily *minimal* with respect to \(g_\infty\), or can persistent **non‑minimal Ricci‑biased interfaces** occur when \(\alpha\neq0\)?  

The conjecture asks for an intrinsic limit theory for curvature‑dependent phase fields on manifolds whose metrics undergo Ricci evolution, including precise identification of any **Ricci‑flow forcing** term that modifies classical mean curvature motion.

Why is it a "good" problem:  
This problem refines the Allen–Cahn–mean curvature correspondence into a dynamic geometric context where the ambient metric itself evolves.  It avoids speculative formulas and instead targets the analytical core: whether a **Brakke‑type compactness and energy‑dissipation principle** continues to hold when both the curvature and the Allen–Cahn field co‑evolve.  
Resolving it demands integrating techniques from geometric measure theory (varifold convergence and Brakke inequalities), parabolic PDEs with time‑dependent metrics, and Ricci flow analysis—areas currently lacking a unified framework.  
Success would produce a new analytic bridge between **Ricci flow entropy methods** and **phase‑field curvature dynamics**, revealing how geometric deformation influences energy concentration and interface motion.  
The statement is conceptually simple but penetrates deep, unexplored territory at the intersection of evolving geometry, Γ‑convergence, and curvature‑driven pattern formation—qualities that mark it as a genuinely “good’’ mathematical problem.